\documentclass{plantillaPPT}

\titulo{9}{Repaso y Parametrización}

\begin{document}

\primeraDiapo

\begin{frame}{Sólidos de Revolución}
\framesubtitle{Práctica 9}

2. Sea \(R\) la región acotada por las curvas.

\begin{center}
\(C_1: y=x+3 \qquad C_2:y=x^3-3 \qquad C_3: y=5-x^2\)
    \includegraphics[width=0.55\textwidth]{C2 2025-1/imagenes/ayud9img.png}
\end{center}

\end{frame}

\begin{frame}{Sólido de Revolución}
\framesubtitle{Práctica 9}

\((a)\) Utilizando el método de capas cilíndricas, escriba la integral que permite calcular el volumen del sólido de revolución que se genera al hacer rotar la región \(R\) en torno a la recta \(x=3\). \\[5ex]
\((b)\) Utilizando el método de discos, escriba la integral que permite calcular el volumen del sólido de revolución que se genera al hacer rotar la región \(R\) en torno a la recta \(y=6\).

    
\end{frame}

\begin{frame}{Sólidos de Revolución}
\framesubtitle{Práctica 9}

\((c)\) Escribir la integral que permite calcular el perímetro de la región \(R\).\\[5
ex]
\((d)\) Escribir la integral que permite calcular el área de la superficie de revolución que se genera al hacer rotar la región \(R\) en torno a la recta \(x=-3\).

\end{frame}

\begin{frame}{Parametrización}
\framesubtitle{Práctica 9}

2. Determine una parametrización para la curva \(C\), que se indica en cada caso. Además, escriba las integrales que permiten calcular la longitud de arco de cada una de ellas.\\[4ex]
\((a)\) \(C\) es el segmento orientado que une los puntos \(A(1, 2)\) y \(B(4, 1)\).\\[2ex]

\((b)\) \(C\) es el arco de la circunferencia de ecuación \(x^2+y^2=4\) y que va desde el punto \(A(2, 0)\) hasta \(B(-\sqrt{3}, 1)\), con \(y\geq0\).

\end{frame}

\begin{frame}{Uso de Paramétricas}
\framesubtitle{Práctica 9}

3. Determine los valores de \(p\in \mathbb{R}\) para que la siguiente integral impropia:
\[\int_{0}^{\infty}\frac{1}{(x+1)^p}\,dx\]

es convergente.



\end{frame}

\end{document}